\documentclass[11pt,a4paper]{article}
\usepackage[utf8]{inputenc}
\usepackage{geometry}
\geometry{margin=1in}
\usepackage{hyperref}
\usepackage{enumitem}
\usepackage{listings}
\usepackage{xcolor}
\usepackage{booktabs}

\lstset{
  language=C,
  basicstyle=\ttfamily\small,
  keywordstyle=\bfseries,
  commentstyle=\itshape,
  showstringspaces=false,
  breaklines=true,
  captionpos=b,
  numbers=left,
  numberstyle=\tiny\color{gray},
}

\title{Digital Design Training Program \\ \large Day 5: The Bare-Metal C Transition}
\author{Authored by }
\date{April 1, 2026}

\begin{document}

\maketitle

\section*{Theme: Stripping away the Arduino framework to understand the underlying AVR hardware architecture, memory registers, and bitwise mathematics.}

\subsection*{Introduction: Why Go Bare-Metal?}
The Arduino framework (\verb|digitalWrite|, \verb|pinMode|, etc.) is fantastic for getting started quickly. However, it hides what the microcontroller is actually doing. By writing ``bare-metal'' C code, we speak directly to the hardware registers. The benefits are immense:
\begin{itemize}
    \item \textbf{Speed:} A direct register manipulation can be hundreds of times faster than a \verb|digitalWrite()| call.
    \item \textbf{Size:} Bare-metal programs are significantly smaller, saving precious flash memory and RAM.
    \item \textbf{Control:} You gain access to the full power of the microcontroller, including features not exposed by the Arduino API.
    \item \textbf{Understanding:} You will finally understand \emph{how} \verb|digitalWrite| actually works.
\end{itemize}
From now on, we will be writing in pure C (\verb|.c| files) and using the standard \verb|main()| function as our entry point.

\subsection*{Module 5.1: Bitwise Operations and Masking (45 Minutes)}

\begin{itemize}
    \item \textbf{Goal:} Master the bitwise operators, the fundamental tools for manipulating hardware registers.
    \item \textbf{Conceptual Overview:} A hardware register is just a byte (8 bits) of memory. Each bit often controls a specific feature. We cannot simply write a new value to the register, as that would change all 8 bits at once. We need to surgically modify only the specific bit we care about.
    \item \textbf{The Operators and Their Use Cases:}
    \begin{enumerate}[label=\arabic*.]
    \item \textbf{Setting a Bit (Bitwise OR} \verb#|#): Use this to turn a feature ON.
        \begin{itemize}
            \item \textbf{Scenario:} A register \verb|REG| is \verb|0b00001000|. We want to turn on bit 3 (the 4th bit from the right) without touching the others.
            \item \textbf{Method:} We create a ``mask'' with a 1 in the target position: \verb|(1 << 3)|, which is \verb|0b00001000|. Then we OR it with the original value.
            \item \textbf{Code:} \verb|REG |= (1 << 3);|
            \item \textbf{Math:} \verb|0b00001000 | 0b00001000| results in \verb|0b00011000|. Bit 3 is now ON.
        \end{itemize}
\item \textbf{Clearing a Bit (Bitwise AND} \verb|&| \textbf{and NOT} \verb|~|): Use this to turn a feature OFF.
        \begin{itemize}
            \item \textbf{Scenario:} A register \verb|REG| is \verb|0b00011000|. We want to turn off bit 3.
            \item \textbf{Method:} We create the same mask \verb|(1 << 3)|. We then invert it with \verb|~| to get \verb|0b11110111|. This mask has a 0 at the target position and 1s everywhere else. We then AND this with the original value.
            \item \textbf{Code:} \verb|REG &= ~(1 << 3);|
            \item \textbf{Math:} \verb|0b00011000 & 0b11110111| results in \verb|0b00010000|. Bit 3 is now OFF.
        \end{itemize}
\item \textbf{Toggling a Bit (Bitwise XOR} \verb|^|): Use this to flip a bit's state.
        \begin{itemize}
            \item \textbf{Scenario:} We want to make an LED blink without caring if it's currently on or off.
            \item \textbf{Code:} \verb|REG ^= (1 << 3);|
            \item \textbf{Math:} If bit 3 is 0, \verb|0 ^ 1 = 1|. If bit 3 is 1, \verb|1 ^ 1 = 0|. It flips automatically.
        \end{itemize}
\item \textbf{Checking a Bit (Bitwise AND} \verb|&|): Use this to read the state of a feature.
        \begin{itemize}
            \item \textbf{Scenario:} We want to know if a button connected to bit 3 is pressed.
            \item \textbf{Code:} \verb|if (REG & (1 << 3)) { ... }|
            \item \textbf{Math:} \verb|0b00011000 & 0b00001000| results in \verb|0b00001000|, which is non-zero (evaluates to \verb|true|). If bit 3 were 0, the result would be \verb|0b00000000| (evaluates to \verb|false|).
        \end{itemize}
    \end{enumerate}
    \item \textbf{Checkpoint:} You can confidently write the one-line C expression to set, clear, toggle, or check any given bit in a byte without affecting the other bits.
\end{itemize}
\newpage
\subsection*{Module 5.2: AVR Architecture and Memory Maps (45 Minutes)}

\begin{itemize}
    \item \textbf{Goal:} Understand how physical microcontroller pins map to specific memory addresses and the three key registers that control them.
    \item \textbf{Requirements:} ATmega328P Pinout Diagram and Datasheet.
    \item \textbf{Conceptual Overview:} The name ``Arduino Pin 8'' is an abstraction. The ATmega328P microcontroller at the heart of the Arduino doesn't know this name. It knows its pins by ``Ports''. Most pins are grouped into 8-bit Ports, named \verb|PORTB|, \verb|PORTC|, and \verb|PORTD|. ``Arduino Pin 8'' is physically \verb|Port B, Bit 0|, which we call \verb|PB0|.
    \item \textbf{The Three Essential Registers:} For each Port (e.g., Port B), there are three 8-bit registers in memory that control its behavior.
    \begin{itemize}
        \item \textbf{\texttt{DDRB} (Data Direction Register for Port B):} This register controls whether the 8 pins of Port B are inputs or outputs.
            \begin{itemize}
                \item To make a pin an \textbf{OUTPUT}, you set its corresponding bit in the DDR to \textbf{1}.
                \item To make a pin an \textbf{INPUT}, you clear its corresponding bit in the DDR to \textbf{0}.
            \end{itemize}
        \item \textbf{\texttt{PORTB} (Data Register for Port B):} This register has a dual purpose.
            \begin{itemize}
                \item If a pin is configured as an \textbf{OUTPUT}, setting its bit in \verb|PORTB| to \textbf{1} drives the pin HIGH (5V). Clearing it to \textbf{0} drives it LOW (GND).
                \item If a pin is configured as an \textbf{INPUT}, setting its bit in \verb|PORTB| to \textbf{1} enables the internal pull-up resistor. Clearing it to \textbf{0} disables the pull-up.
            \end{itemize}
        \item \textbf{\texttt{PINB} (Port Input Pins Register for Port B):} This is a read-only register. Reading from it gives you the current logical state (1 or 0) of all 8 pins in the port.
    \end{itemize}
    \item \textbf{Arduino to AVR Translation Table:}
    \begin{center}
    \begin{tabular}{ll}
    \toprule
    \textbf{Arduino Code} & \textbf{Bare-Metal C Equivalent (for Arduino Pin 8 / PB0)} \\
    \midrule
    \verb|pinMode(8, OUTPUT);| & \verb|DDRB |= (1 << PB0);| \\
    \verb|pinMode(8, INPUT);| & \verb|DDRB &= ~(1 << PB0);| \\
    \verb|pinMode(8, INPUT_PULLUP);| & \verb|DDRB &= ~(1 << PB0); PORTB |= (1 << PB0);| \\
    \verb|digitalWrite(8, HIGH);| & \verb|PORTB |= (1 << PB0);| \\
    \verb|digitalWrite(8, LOW);| & \verb|PORTB &= ~(1 << PB0);| \\
    \verb|digitalRead(8);| & \verb|(PINB & (1 << PB0))| \\
    \bottomrule
    \end{tabular}
    \end{center}
    \item \textbf{Checkpoint:} Given an Arduino pin number, you can use the pinout diagram to find its \verb|Pxy| name and then write the correct C expressions to configure its direction, set its state, or read its value.
\end{itemize}
\newpage
\subsection*{Module 5.3: Bare-Metal I/O Translation (60 Minutes)}

\begin{itemize}
    \item \textbf{Goal:} Rewrite Day 1's simple LED and Button logic using strictly bare-metal C register manipulation.
    \item \textbf{Action Steps \& Code:} We will translate the simple button-and-LED sketch from Day 1 into pure C. The LED is on Arduino Pin 8 (\verb|PB0|) and the button is on Arduino Pin 9 (\verb|PB1|).
\begin{lstlisting}[caption={Code for Module 5.3: Bare-Metal Blink}]
// avr/io.h contains all the definitions for the registers like PORTB, DDRB, etc.
#include <avr/io.h>

// The program entry point is now main(). It returns an integer and takes no arguments.
int main(void) {
    // --- This is our ``setup()'' ---

    // Set Arduino Pin 8 (PB0) as an output.
    // Equivalent to: pinMode(8, OUTPUT);
    DDRB |= (1 << PB0);

    // Set Arduino Pin 9 (PB1) as an input.
    // Equivalent to: pinMode(9, INPUT);
    DDRB &= ~(1 << PB1);

    // Enable the pull-up resistor on Pin 9 (PB1).
    // Equivalent to: digitalWrite(9, HIGH) on an input, or the INPUT_PULLUP mode.
    PORTB |= (1 << PB1);

    // --- This is our ``loop()'' ---
    // The while(1) loop runs forever.
    while (1) {
        // Check if Pin 9 (PB1) is LOW.
        // The PINB register gives the state of the physical pins.
        // We check if the PB1 bit is 0.
        // Equivalent to: if (digitalRead(9) == LOW)
        if (!(PINB & (1 << PB1))) {
            // If button is pressed, turn LED ON.
            // Drive Pin 8 (PB0) HIGH.
            // Equivalent to: digitalWrite(8, HIGH);
            PORTB |= (1 << PB0);
        } else {
            // If button is not pressed, turn LED OFF.
            // Drive Pin 8 (PB0) LOW.
            // Equivalent to: digitalWrite(8, LOW);
            PORTB &= ~(1 << PB0);
        }
    }
    return 0; // This line is never reached.
}
\end{lstlisting}
    \item \textbf{Checkpoint:} The breadboard setup from Day 1 (button controls LED) functions identically. If you compile this code and the original Arduino sketch, you will see that this bare-metal version uses significantly less program storage space.
    \item \textbf{Challenge:} Modify the \verb|while(1)| loop. Instead of the \verb|if/else| block, can you make the LED toggle every time the button is pressed? This will require combining the debouncing logic from Day 4 with the bitwise \verb|XOR| operator (\verb|^|).
\end{itemize}

\end{document}
